\paragraph{Choosing $F_\text{obs}$ in practice.}
The three constructors have different epistemic credentials. Dataset quality bounds (Constructor~1) are the most defensible because they are external to the forecaster being evaluated; the low 9.9\% reversal rate under Constructor~1 (\S5) itself suggests that UCDP coding ranges understate true observational uncertainty. Parametric defaults (Constructor~3) are practical but introduce a free parameter; a careless choice of $\sigma_\text{rel}$ can mask or invert the ranking, and benchmarking communities would need to agree on standard values per data type. The RVI Gumbel mixture (Constructor~2) bridges these approaches: it uses regression coefficients estimated from coder surveys \citep{vesco2026uncertainty} to predict the distribution of plausible fatality counts, capturing the asymmetric underreporting that symmetric parametric constructors miss. The Vesco et al.\ model has its own limitations---response biases in coder surveys, assumed distributional form, limited sample size---and should be read as a principled estimate rather than ground truth. That said, I recommend Constructor~2 where UCDP data are available, Constructor~1 where quality bounds exist for other datasets, and Constructor~3 with a community-agreed $\sigma_\text{rel}$ otherwise. A caveat: the Constructor~2 results in \S5 apply the RVI regression at the cell-month aggregate level rather than per event, introducing a unit-of-analysis mismatch that likely understates the true observational uncertainty (see \S5); event-level application would be more principled but requires combining event-level distributions. A fourth possibility---deriving $F_\text{obs}$ from a model's own confidence signal---is conceptually available but introduces circularity; I do not pursue it here.

\paragraph{Relation to prior approaches.}
The problem of evaluating forecasts against imprecise observations is well studied; the present work applies ideas developed primarily in meteorology to conflict forecasting. One structural consequence worth noting first: \citet{brocker2009reliability} decomposes proper scores into reliability, resolution, and uncertainty components; the Cram\'er-distance reformulation shifts mass from the ``uncertainty'' term---which in the classical decomposition absorbs all observational noise---into the smooth $F_\text{obs}$ in the integrand. A formal restatement of this decomposition for the reformulated metric is deferred to a technical companion.

\textbf{Scoring under observation error.} \citet{brocker2007scoring} highlight the need for generalised scores when verification data are uncertain, proposing a convolution approach: the forecast is convolved with the observation-error kernel before scoring, so that forecast and observation are compared on equal footing. \citet{bessac2021forecast} take a different route, computing $\mathbb{E}[S(F, Y) \,|\, \tilde{y}]$---the conditional expectation of the score over the observation error distribution---yielding an unbiased estimate of the true score. This conditional-expectation correction is principled when the observation error model is known. The Cram\'er-distance reformulation takes a third route: rather than correcting the score after the fact, it replaces the degenerate observation CDF with a smooth $F_\text{obs}$ \emph{inside the integrand}, changing what is being measured. The Br\"ocker--Smith convolution asks ``what would the score look like if the forecast absorbed the observation noise?''; the Bessac--Naveau correction asks ``what would the score have been had we observed the truth?''; the Cram\'er-distance reformulation asks ``what is the distance between the forecast distribution and the observation distribution?'' The three approaches yield the same ranking in expectation when $F_\text{obs}$ is correctly specified, but they differ in what they require: the convolution approach requires a noise kernel, the conditional-expectation approach requires a tractable posterior $[X \mid Y]$, and the CDF-substitution approach requires a constructable $F_\text{obs}$.

I verify this empirically by comparing the Cram\'er-distance reformulation against the Bessac--Naveau conditional-expectation correction. On the controlled demonstration (\S5), both approaches select the same winner (honest version) for all $\sigma_\text{honest} \leq 0.80$ and agree that the tight version wins for $\sigma_\text{honest} \geq 0.90$. A narrow disagreement band exists at $\sigma_\text{honest} \approx 0.85$--$0.88$, where the Cram\'er distance still selects the honest version but the conditional-expectation correction flips to tight---the CDF-substitution approach is slightly more tolerant of wide versions near the transition. On the UCDP empirical comparison (19{,}397 cell-months, Constructor~3 with $\sigma_\text{rel} = 0.4$), the two approaches agree in 98.1\% of cell-months overall. Outside the transition zone ($\hat{\sigma}_\text{cell} \leq 0.8$ or $> 1.0$) concordance is 100\%; within the transition zone ($0.8 < \hat{\sigma}_\text{cell} \leq 1.0$) it drops to 93.6\%---roughly one in fifteen cell-months receives a different ranking depending on the method (\S5). The residual disagreement is concentrated precisely where the decision is hardest: where observational uncertainty is large enough to be consequential but not so large as to swamp the forecast signal. In this disagreement band, the conditional-expectation correction has the stronger theoretical warrant---it recovers an unbiased estimate of the true score under correct error-model specification---whereas the Cram\'er distance measures a different quantity (distribution-to-distribution distance) whose optimality depends on the $F_\text{obs}$ constructor.

\textbf{Error-corrected and distributional approaches.} \citet{ferro2017observation} constructs error-corrected proper scoring rules whose expected values, taken over the observation-error distribution, equal the original proper score evaluated at the true value---a frequentist unbiasedness condition that requires an explicit error model but not a Bayesian posterior. \citet{candille2008impact} study the impact of observation error on ensemble verification diagnostics, proposing an ``observational-probability'' method that replaces point observations with probability distributions in scores such as the Brier score and reliability diagrams. \citet{bessac2021forecast} derive the full distribution of CRPS under Gaussian additive and multiplicative-Gamma observation errors, showing that imperfect verification data cause systematic underdispersion bias---an effect closely related to the ranking reversal demonstrated in \S5.

The underdispersion bias Bessac and Naveau predict analytically is directly observable in the UCDP data (\S5): classical CRPS selects the tight version in 100\% of cell-months, and the Cram\'er-distance reformulation reverses this in precisely the regime where observational uncertainty exceeds forecast uncertainty (Spearman $\rho = -0.38$, $p < 10^{-15}$, between $\sigma_\text{honest}$ and the classical score gap).

These approaches share the same motivation as the present work but address different aspects: score bias (Bessac--Naveau), error-corrected scoring (Ferro), or diagnostic calibration (Candille--Talagrand), whereas I focus on the effect on model ranking.

\textbf{Scale-sensitivity and ensemble size.} \citet{bolin2023local} show that CRPS is not locally scale-invariant: it overweights observations in regions where the forecast distribution has large variability. They propose a Scaled CRPS that normalises by local spread---a complementary approach that could in principle be combined with the reformulation proposed here. \citet{ferro2017fair} addresses finite ensemble size bias, a distinct but related source of overconfidence in CRPS rankings. \citet{weijs2011accounting} cast the problem in information-theoretic terms, proposing divergence-based scores with explicit uncertainty decompositions; their framework is more general but requires different scoring infrastructure.

Among these approaches, the Cram\'er-distance reformulation is attractive for domains where observation error models are not well characterised: it requires only a constructable $F_\text{obs}$, not an invertible error model, and reduces exactly to classical CRPS when observations are precise. Its main advantage for conflict forecasting is that it provides a natural entry point for empirically grounded uncertainty models: UCDP's own low/best/high bounds are coding artefacts that understate the true observational uncertainty (\S5), but the RVI Gumbel mixture of \citet{vesco2026uncertainty}---built precisely because those bounds are insufficient---can be plugged directly into the $F_\text{obs}$ constructor. Fields with well-characterised instrument error covariances (e.g.\ numerical weather prediction) may benefit more from the conditional-expectation correction of \citet{bessac2021forecast}, the convolution approach of \citet{brocker2007scoring}, or the error-corrected scoring rules of \citet{ferro2017observation}.

\paragraph{What the reformulation does not solve.}
\begin{itemize}
    \item \textbf{It does not solve the Forecaster's Dilemma.} The result of \citet{lerch2017forecasters}---that conditioning forecast evaluation on extreme observations systematically discredits skillful forecasts---holds for the reformulated metric just as for classical CRPS. Weighted scoring rules that emphasise tail behaviour \citep{gneiting2011comparing} are still needed as complementary diagnostics.

    \item \textbf{It does not eliminate distributional assumptions.} The constructor of $F_\text{obs}$ is itself a distributional assumption. The reformulation moves the assumption from ``observations are exact'' to ``observations follow this distribution''---more honest, but still an assumption open to mis-specification.

    \item \textbf{Additional sources of observational uncertainty are unaddressed.} Systematic data-quality biases, spatial variation in reporting quality, georeferencing imprecision (UCDP precision codes 1--7), and false-negative uncertainty in the roughly 800{,}000 zero-fatality cell-months excluded from \S5 each represent distinct error sources not captured by $F_\text{obs}$.

    \item \textbf{The demonstrations use oracle-centred forecast versions.} The controlled and UCDP comparisons centre both versions on the reported value $y$---information unavailable at prediction time. The 39\% reversal rate describes a regime of perfect location knowledge; the HydraNet validation (\S5.3) tests whether the reversal survives genuine location error. All three comparisons test uncertainty valuation rather than model selection---whether the reformulation changes rankings between genuinely different models is untested.
\end{itemize}

\paragraph{Sensitivity to $F_\text{obs}$ mis-specification.}
Could the metric's behaviour under a mis-specified $F_\text{obs}$ be worse than classical CRPS with no $F_\text{obs}$ at all? I quantify this by running the stochastic consistency check from \S4 (details in Appendix~\ref{sec:appendix_consistency}) under deliberate mis-specification. The true data-generating distribution is $\text{LogNormal}(\mu = 1.0, \sigma = 0.3)$; I construct $F_\text{obs}$ using $\sigma_\text{rel}$ values ranging from $0.1$ to $1.0$ and ask where the expected Cram\'er distance is minimised.

Table~\ref{tab:misspec} shows the results. When $\sigma_\text{rel}$ matches the true $\sigma$ (ratio 1.0$\times$), the argmin falls exactly on $\mu_\text{true}$. Moderate overestimation (up to 2$\times$) shifts the argmin by at most one grid step ($+0.05$). Severe overestimation (3.3$\times$) shifts it by $+0.15$---meaningful but bounded. Underestimation produces no detectable bias in this sweep. For this log-normal DGP, the bias is monotonic in the mis-specification ratio; whether monotonicity holds across distributional families is untested, but the practical implication---a practitioner can bound the bias by bounding the mis-specification---applies whenever the relationship is monotone.

The asymmetry has a structural explanation. When $\sigma_\text{rel}$ is underestimated, $F_\text{obs}$ approaches the degenerate step function at $y$, and the Cram\'er distance reduces to classical CRPS---which is proper, so the location argmin remains at $\mu_\text{true}$. When $\sigma_\text{rel}$ is overestimated, $F_\text{obs}$ is wider than the true observation noise; the log-normal right-skew of the widened $F_\text{obs}$ shifts the $L^2$-minimising location upward, producing positive bias. The same mechanism explains the scale-consistency convolution effect discussed in \S4 and Appendix~\ref{sec:appendix_consistency}: the metric measures distance from the observable distribution, which inherits the observation noise.

A formal characterisation of this bias is an open problem.

\begin{table}[h]
    \centering
    \small
    \begin{tabular}{rrrr}
        \toprule
        $\sigma_\text{rel}$ & Ratio ($\sigma_\text{rel} / \sigma_\text{true}$) & Argmin $\hat{\mu}$ & Bias ($\hat{\mu} - \mu_\text{true}$) \\
        \midrule
        0.1 & 0.3$\times$ & 1.000 & 0.000 \\
        0.2 & 0.7$\times$ & 1.000 & 0.000 \\
        0.3 & 1.0$\times$ & 1.000 & 0.000 \\
        0.4 & 1.3$\times$ & 1.050 & $+$0.050 \\
        0.6 & 2.0$\times$ & 1.050 & $+$0.050 \\
        0.8 & 2.7$\times$ & 1.100 & $+$0.100 \\
        1.0 & 3.3$\times$ & 1.150 & $+$0.150 \\
        \bottomrule
    \end{tabular}
    \caption{Argmin bias under mis-specified $F_\text{obs}$. The true DGP is $\text{LogNormal}(1.0, 0.3)$; $F_\text{obs}$ uses varying $\sigma_\text{rel}$. Monte Carlo expectation over 500 draws. Grid step $= 0.05$.}
    \label{tab:misspec}
\end{table}

\paragraph{Implications for forecast competitions.}
Forecast competitions and model selection procedures that use CRPS as the sole ranking metric implicitly assume observations are exact. When this assumption is violated, competition rankings can systematically favour overconfident models. The practical recommendation is not to replace CRPS with the Cram\'er-distance reformulation in all settings---classical CRPS is exactly right when observations are precise---but to use the reformulation as a secondary metric or robustness check in domains where observational uncertainty is known to be substantial. This recommendation rests on location-consistency verified numerically across three distributional families (\S5; Appendix~\ref{sec:appendix_demo}, Table~\ref{tab:multi_dgp}). Domains where this is immediately relevant include conflict forecasting \citep{hegre2019views,ucdp_codebook}, epidemiological surveillance \citep{gibbons2014measuring}, and economic forecasting from survey data.

\paragraph{Open questions.}
\begin{itemize}
    \item Closed-form expressions for the reformulated metric when both $F_\text{pred}$ and $F_\text{obs}$ belong to common parametric families (log-normal, gamma, Weibull) would speed evaluation and enable analytical optimisation.
    \item A formal Br\"ocker decomposition of the reformulated metric into reliability, resolution, and observational-uncertainty components.
    \item An analytical proof of consistency under stochastic observations via the energy-distance framework, with explicit conditions on the $F_\text{obs}$ constructor.
    \item Empirical comparison of model rankings on operational benchmarks (conflict forecasting, weather prediction, epidemiological surveillance) under classical and reformulated CRPS with each constructor.
    \item Convergence analysis: how does the bias from a mis-specified $F_\text{obs}$ scale with the magnitude and direction of the mis-specification?
\end{itemize}
